% Image credits for the photographs and AI illustrations of Book 2 (High School).
% Machine-readable license records live in images/book2/CREDITS.md; keep the
% two files in sync. The Book 2 writer owns this file: add one \item per
% photograph (inside the itemize, by chapter) as photographs land.
\chapter*{Image Credits}

The drawings, maps and charts in this book are original; maps and charts
drawn from published data name their source in the caption or in this
list. Photographs are used under their respective free licenses, with
thanks to their authors; full source links and license URLs for every
photograph are listed in \texttt{images/book2/CREDITS.md} in the book's
repository.

\bigskip
Alongside the photographs and the original drawings, some figures are
AI-generated illustrations, produced with OpenAI image generation from
prompts written by the book's editors, and reviewed for scientific
accuracy before inclusion. They show generic scenes only, never a real
place, a real person or a real specimen. The full prompt for every
generated image is recorded in \texttt{images/book2/ai/PROMPTS.md} in the
repository.

\bigskip
\noindent Photographs and data, by chapter:
\begin{itemize}\small
  \item Ch.~1 (measuring the Earth): the Earth seen from Apollo~17; photograph by NASA, public domain (Wikimedia Commons).
  \item Ch.~1: a repeating circle, plate from J.-B. Delambre, \emph{Base du système métrique décimal}, vol.~2 (1807); public domain (Wikimedia Commons, from BnF Gallica).
  \item Ch.~1: degree lengths computed on the WGS~84 ellipsoid (EPSG dataset).
  \item Ch.~2 (a habitable planet): Venus from Mariner~10 (PIA23791), NASA/JPL-Caltech, public domain; Mount Sharp from the rover Curiosity (PIA16769), NASA/JPL-Caltech/MSSS, public domain (Wikimedia Commons).
  \item Ch.~2: planetary data from the NASA fact sheets; water phase boundaries from the IAPWS releases SR1-86 and R14-08; habitable-zone limits from Kopparapu and others (2013); Proxima~b from the NASA Exoplanet Archive.
  \item Ch.~3 (the Earth's energy budget): melt ponds on Arctic sea ice (ICESCAPE mission, 2011); photograph by NASA Goddard Space Flight Center, public domain (Wikimedia Commons).
  \item Ch.~3: energy-budget values from the IPCC Sixth Assessment Report, Working Group~I, Figure~7.2; albedos from the University of Colorado Boulder (MOSAiC) and the National Snow and Ice Data Center; heat transport from Trenberth and Caron (2001); net primary production from Field and others (1998); thermochemical data from the NIST Chemistry WebBook.
  \item Ch.~4 (why clouds form): fair-weather cumulus, photograph by Ralph F. Kresge, NOAA; the anvil of a cumulonimbus, NOAA National Severe Storms Laboratory; both public domain (Wikimedia Commons).
  \item Ch.~4: saturation computed from the IAPWS releases; drop sizes from the NASA Jet Propulsion Laboratory.
  \item Ch.~5 (storms and tropical cyclones): Hurricane Isabel from the International Space Station, photograph by Mike Trenchard, NASA Johnson Space Center; a tornado in Kansas, photograph by Sean Waugh, NOAA National Severe Storms Laboratory; lightning, photograph by C.~Clark, NOAA National Severe Storms Laboratory; all public domain (Wikimedia Commons).
  \item Ch.~5: cyclone tracks from IBTrACS version~4r01 (Gahtan and others, 2024; Knapp and others, 2010; NOAA National Centers for Environmental Information); hurricane categories from the NOAA National Hurricane Center; hurricane energy and formation conditions from the NOAA Atlantic Oceanographic and Meteorological Laboratory; flash rate from Christian and others (2003); tornado facts from the NOAA Storm Prediction Center.
  \item Ch.~6 (the layered ocean): a sampling rosette, photograph by Pablo Clemente-Colon, NOAA Photo Library; a profiling float launched from a research ship, NOAA Global Ocean Monitoring and Observing Program; ocean chlorophyll 2002--2010, Jesse Allen and Robert Simmon, NASA Earth Observatory; all public domain (Wikimedia Commons).
  \item Ch.~6: ocean profiles from the World Ocean Atlas 2023 (Reagan and others, 2024; NOAA National Centers for Environmental Information); seawater density from the international equation of state (Fofonoff and Millard, 1983); ocean volume from Charette and Smith (2010); float facts from the Argo Program Office; upwelling from the NOAA National Ocean Service and Kessler (2006).
  \item Ch.~7 (soil, a slow-forming resource): a laterite profile on gneiss in Kerala, photograph by Werner Schellmann, CC~BY-SA~2.5; saprolite in central Madagascar, photograph by Diorit, CC~BY-SA~3.0; a soil profile, USDA Natural Resources Conservation Service, public domain (all Wikimedia Commons).
  \item Ch.~7: soil production and erosion rates from Montgomery (2007); mineral formulas from the Mineralogy Database (webmineral.com); saprolite thickness from the USGS Ground Water Atlas of the United States.
  \item Ch.~8 (drainage basins, rivers and floods): a hydrologist measuring streamflow and a gauging station, U.S. Geological Survey, public domain (Wikimedia Commons).
  \item Ch.~8: annual peaks and daily discharge of the Potomac River near Washington, U.S. Geological Survey National Water Information System, station 01646500; Mississippi discharge and basin from Kammerer (1990, USGS); Amazon discharge from Call\`ede and others (2010); flood-water safety from the NOAA National Weather Service.
  \item Ch.~9 (freshwater resources): the Aral Sea in 2000 and 2018, NASA Earth Observatory (MODIS/Terra); irrigated fields near Sublette, Kansas, NASA (ASTER/Terra), public domain (Wikimedia Commons).
  \item Ch.~9: world water stocks after Shiklomanov (1993), U.S. Geological Survey Water Science School; the global water cycle after Trenberth and others (2007), as printed in Trenberth and Asrar (2014); Nubian groundwater ages from Sturchio and others (2004); High Plains aquifer changes from McGuire (2017, USGS); the Aral Sea budget from Gaybullaev and others (2012) and Micklin (2007).
  \item Ch.~10 (from continental drift to seafloor spreading): the drilling ship \emph{JOIDES Resolution}, Ocean Drilling Program, Texas A\&M University, public domain (Wikimedia Commons).
  \item Ch.~10: the polarity time scale of Cande and Kent (1995); the Leg~3 drilling results of the Deep Sea Drilling Project (1970); the age of the ocean floor from Seton and others (2020), EarthByte, licence CC~BY~4.0.
  \item Ch.~11 (seismic waves reveal the interior): installing a temporary seismometer in Puerto Rico; photograph by Thomas Pratt, U.S. Geological Survey, public domain (Wikimedia Commons).
  \item Ch.~11: the Preliminary Reference Earth Model of Dziewonski and Anderson (1981); empirical travel times of P and S waves from the Australian Passive Seismic Server (AuSPASS).
  \item Ch.~12 (lithosphere and asthenosphere): an iceberg, NOAA, public domain (USGS Water Science School).
  \item Ch.~12: heat flow from Davies and Davies (2010), licence CC~BY~3.0; heat-flow provinces from Roy, Blackwell and Decker (1972); the mantle solidus of Hirschmann (2000); crust thickness and density from Christensen and Mooney (1995); elevations from ETOPO1 (NOAA).
  \item Ch.~13 (plates and their motions): a satellite-positioning station in the Cascade Range; photograph by Dave Sherrod, U.S. Geological Survey, public domain (Wikimedia Commons). Shaded relief of the North Pacific floor, National Geophysical Data Center, NOAA, public domain (Wikimedia Commons).
  \item Ch.~13: plate rotation poles NNR-MORVEL56 of Argus, Gordon and DeMets (2011); station positions from the Nevada Geodetic Laboratory (Blewitt, Hammond and Kreemer, 2018); volcano ages of Clague and Dalrymple (1987), U.S. Geological Survey; plate boundaries of Bird (2003); land from Natural Earth.
  \item Ch.~14 (ocean ridges): pillow lava on the Galápagos Rift, NOAA Okeanos Explorer Program, public domain; a black smoker on the Juan de Fuca Ridge, University of Washington and NOAA, public domain; gabbro, photograph by Darla Sondrol (GeoDIL), CC0; peridotite in basalt, photograph by James St.~John, CC~BY~2.0 (all Wikimedia Commons).
  \item Ch.~14: ocean-floor depths from ETOPO1 (NOAA) and ages from Seton and others (2020), EarthByte, licence CC~BY~4.0; depth laws of Parsons and Sclater (1977); the mantle solidus of Hirschmann (2000); rotation poles NNR-MORVEL56 of Argus, Gordon and DeMets (2011).
  \item Ch.~15 (subduction zones): blueschist and eclogite, photographs by James St.~John, CC~BY~2.0; tsunami damage at Kesennuma, Japan, 2011, photograph by Dylan McCord, U.S. Navy, public domain (all Wikimedia Commons).
  \item Ch.~15: earthquakes from the USGS earthquake catalogue (ComCat); plate boundaries and convergence rates of Bird (2003); ocean-floor ages from Seton and others (2020), EarthByte, licence CC~BY~4.0; plate speeds NNR-MORVEL56 of Argus, Gordon and DeMets (2011); land from Natural Earth.
  \item Ch.~16 (minerals): quartz, orthoclase, biotite and olivine, photographs by Shannon Heinle, Darla Sondrol, Dexter Perkins and Nessa Eull (Geology Digital Image Library), CC0; basalt in thin section, photographs by Chmee2, CC~BY-SA~3.0; quartz in a granite in thin section, photograph by Michael C.~Rygel, CC~BY-SA~3.0 (all Wikimedia Commons).
  \item Ch.~16: the composition of the continental crust of Rudnick and Gao (2014); standard atomic weights of IUPAC/CIAAW; mineral formulas and densities from webmineral.com; densities of Sierra Nevada granitic rocks from the U.S. Geological Survey (Open-File Report 91-570-A).
  \item Ch.~17 (igneous rocks): basalt columns of the Giant's Causeway, photograph by Adam Bishop, CC~BY-SA~3.0; El Capitan, Yosemite, U.S. Geological Survey, CC0; basalt, andesite, rhyolite, diorite, granite, obsidian, a porphyritic rock and two thin sections, photographs by Nessa Eull, Twyla Baker, Joseph H.~Hartman, Mike Davis, Dexter Perkins and Darla Sondrol (Geology Digital Image Library), CC0 (all Wikimedia Commons).
  \item Ch.~17: cooling of dykes and sills after Delaney (1986), U.S. Geological Survey Open-File Report 86-444; magma compositions and the order of crystallisation from S.~A.~Nelson's physical geology notes and S.~Earle, \emph{Physical Geology} (CC~BY~4.0); Bowen's biography from the Carnegie Science history of the Geophysical Laboratory.
  \item Ch.~18 (coal, oil and gas): fossil plants in carbon-rich shale, lignite, bituminous coal and anthracite, photographs by James St.~John, CC~BY~2.0; a lignite bed in Romania, U.S. Geological Survey, public domain (all Wikimedia Commons).
  \item Ch.~18: oil window, coal formation and traps from S.~Earle, \emph{Physical Geology} (CC~BY~4.0); carbon content of coal from the U.S. Energy Information Administration; the petroleum system after Magoon and Schmoker (U.S. Geological Survey, 2000); the rift and salt basins of western Central Africa from U.S. Geological Survey Fact Sheet 2023-3028.
  \item Ch.~19 (the continental crust): the Acasta gneiss, photograph by Pedroalexandrade, CC~BY-SA~3.0; Jack Hills zircons, images by A.~Cavosie and J.~Valley, public domain; gneisses and granites of the Canadian Shield, photograph by Chris Light, CC~BY-SA~4.0 (all Wikimedia Commons).
  \item Ch.~19: crust composition of Rudnick and Gao (2014); crustal thickness of Christensen and Mooney (1995), White, McKenzie and O'Nions (1992) and Karplus and others (2019); extent of the continents of Cogley (1984); ocean-floor ages from Seton and others (2020), EarthByte, licence CC~BY~4.0.
  \item Ch.~20 (building mountains): Mount Everest from the space shuttle Columbia, NASA, public domain; folded rock layers in the Kali Gandaki gorge, photograph by Jean-Marie Hullot, CC~BY~2.0; the Glarus Thrust, photograph by Christian Heine, CC~BY-SA~3.0; pillow lava of the Semail ophiolite, photograph by Lionel Duchoiselle, CC~BY-SA~3.0 (all Wikimedia Commons).
  \item Ch.~20: past positions of India and Lhasa from the plate models of M\"uller and others (2019, 2022), computed with the GPlates Web Service (EarthByte); high ground from ETOPO1; earthquakes from the USGS catalogue; plate boundaries of Bird (2003); land from Natural Earth.
  \item Ch.~21 (metamorphic rocks): slaty cleavage under the microscope, photograph by Geoloshe, CC0; garnet mica schist from Syros, photograph by Graeme Churchard, CC~BY~2.0; kyanite schist and quartz--biotite gneiss, photographs by James St.~John, CC~BY~2.0 (all Wikimedia Commons).
  \item Ch.~21: the stability fields of kyanite, andalusite and sillimanite from Holdaway and Mukhopadhyay (1993).
  \item Ch.~22 (magmatism at subduction zones): the eruption plume of Mount St.~Helens, 18 May 1980, photograph by Donald A. Swanson, USGS, public domain; granodiorite, andesite and the Sierra Nevada, photographs by James St.~John, CC~BY~2.0 (all Wikimedia Commons).
  \item Ch.~22: earthquakes from the USGS catalogue; volcanoes from the Smithsonian Global Volcanism Program; plate boundaries of Bird (2003); land from Natural Earth; the dry peridotite solidus of Hirschmann (2000), with PREM densities.
  \item Ch.~23 (the destruction of mountains): the Appalachian Mountains from the International Space Station, astronaut photograph ISS033-E-22378, NASA, public domain; quartz-pebble conglomerate, photograph by James St.~John, CC~BY~2.0 (both Wikimedia Commons).
  \item Ch.~23: denudation rates of Portenga and Bierman (2011); the Mississippi sediment loads of Meade and Moody (2010); the decay constant of Ahnert (1970); the Varden Hills, the series' invented region.
  \item Ch.~24 (the Earth's internal heat): the Nesjavellir geothermal power station, Iceland, photograph by Gretar \'Ivarsson, edited by Fir0002, public domain (Wikimedia Commons).
  \item Ch.~24: heat flow of Davies and Davies (2010); radiogenic heat from the KamLAND geoneutrino measurement (2011) and the present heat release of the four isotopes; heat production of rocks from Vil\`a and Fern\'andez (2010).
  \item Ch.~25 (radiometric dating and the age of the Earth): William Thomson, Lord Kelvin, 1906, photograph by Russell and Sons, London, Smithsonian Institution Libraries, no known copyright restrictions; the Canyon Diablo iron meteorite, photograph by Geoffrey Notkin, CC~BY-SA~2.5 (both Wikimedia Commons).
  \item Ch.~25: meteorite lead isotopes of Patterson (1956); the decay constant of rubidium-87 recommended by IUGS and IUPAC (2022); the history of the estimates after Dalrymple, as summarised by Wikipedia; the Varden Hills, the series' invented region.
  \item Ch.~26 (the early Earth and its atmosphere): banded iron formation, Dales Gorge, photograph by Graeme Churchard, CC~BY~2.0; stromatolites in Shark Bay, photograph by Paul Harrison, CC~BY-SA~3.0; red beds near Sedona, photograph by Wendy from Pennsylvania, CC~BY-SA~2.0 (all Wikimedia Commons).
  \item Ch.~26: the young Sun's luminosity of Gough (1981); eon limits and colours from the International Chronostratigraphic Chart.
  \item Ch.~27 (reading past climates): air bubbles in Antarctic ice, photograph CSIRO Atmospheric Research, CC~BY~3.0; planktic foraminifera, photograph by Marci Robinson, USGS, public domain; pollen grains, Dartmouth Electron Microscope Facility, public domain (all Wikimedia Commons).
  \item Ch.~27: EPICA Dome~C temperature (Jouzel and others, 2007) and the Antarctic ice-core carbon dioxide (Bereiter and others, 2015) via NOAA Paleoclimatology; the LR04 stack (Lisiecki and Raymo, 2005); ETOPO1 relief (NOAA); NASA GISS (GISTEMP v4); reference periods from the IPCC Sixth Assessment Report, Working Group~I, chapter~2.
  \item Ch.~28 (climate change): the Mauna Loa Observatory; photograph by Brian Vasel, NOAA, public domain (Wikimedia Commons).
  \item Ch.~28: Charles David Keeling and John Chin at Mauna Loa, 1997; photograph by John Miller, NOAA, public domain (Wikimedia Commons).
  \item Ch.~28: data from NOAA Global Monitoring Laboratory and Scripps Institution of Oceanography (carbon dioxide), NASA GISS (GISTEMP v4), NOAA Laboratory for Satellite Altimetry (sea level), Bereiter and others (2015) via NOAA Paleoclimatology (ice cores), and the IPCC Sixth Assessment Report, Working Group~I (forcings, attribution, projections).
\end{itemize}
